\section{A fixed algebraic compiler for the twelve-root C1 Gram}
\label{sec:fixed-principal}
This section removes the adaptive testing radicals from the principal
$\lambda>0$ C1 Gram on one explicit testing box. The source admission
condition is finite and algebraic. It is a sufficient condition, and no
claim that it covers every stationary source is made.

Let $\delta=2^{-12}$, $s=\sqrt u$, and write the complete principal Gram
as
\[
 Q(c;s,w)=C_{00}(c)+\sum_{(p,q)\in I} C_{pq}(c)s^pw^q,
 \qquad
 I=\{(0,1),(0,2),(1,0),(2,0),(2,1),(4,0),(4,1)\}.
\]
Here $c$ denotes the source, rather than the angular compensation
scalar; the printed source matrices retain both independent state
normals. Their order is twelve. Define the source reference matrix
$B(c)$ by the diagonal entries
\[
 (2,\,2,\,b^2+(1-b)^2,\,1,\,(K_A+c_A)/2,\,(K_A+c_A)/2,
       1/2,1/2,1/2,1/2,1,1/2)
\]
and the single off-diagonal entry $B_{23}=B_{32}=2b-1$ in one-based
indexing, where $c_A=\max(0,1-K_A)$. Its radial block has determinant
one and trace at most three for $0\le b\le1$. Thus
\[
 B(c)\succeq\tfrac13I_{12}.
\]
An elementary exact verification of its radial reserve is
\[
 (N_y+bR)^2+(N_y-(1-b)R)^2
 =\tfrac13(N_y^2+R^2)
 +\tfrac53\left(N_y+\tfrac35(2b-1)R\right)^2
 +\tfrac{1+6b(1-b)}{15}R^2.
\]
No source-dependent real whitening is used.

For a symmetric source matrix $M$, put
\[
 n(M)=\max_i\sum_j|M_{ij}|.
\]
Every operation here is in the source real closed field, and absolute
values and the finite maximum are semialgebraic functions. Set
\[
 E_0=C_{00}-B,\qquad
 A_{pq}=\delta^{p+q}C_{pq},\qquad r_{pq}=n(A_{pq}).
\]
Our explicit source admission condition is
\begin{equation}\label{eq:finite-admission}
 n(E_0)+\sum_{(p,q)\in I}r_{pq}\le\tfrac16.
\end{equation}
It is imposed together with the printed rational stationary-source
formulas, native positive denominators, $D,\beta\ge0$, and $\lambda>0$.
The condition contains no unknown testing value, no choice of a real
standard part, and no unspecified sufficiently small constant.

The elementary inequality
$2|x_ix_j|\le x_i^2+x_j^2$ proves
$-n(M)I\preceq M\preceq n(M)I$ over every ordered real closed field.
Consequently all the following are positive semidefinite source matrices:
\[
 H_0=C_{00}-\sum_I r_{pq}I\succeq\tfrac16I,
 \qquad H_{pq}^{\pm}=\tfrac12(r_{pq}I\pm A_{pq})\succeq0.
\]
On the fixed testing rectangle $0\le s,w\le\delta$, let
$a=s/\delta$, $b_t=w/\delta$. Direct expansion gives
\begin{equation}\label{eq:fixed-matrix-split}
 Q=H_0+\sum_I\bigl[H_{pq}^+(1+a^pb_t^q)
                         +H_{pq}^-(1-a^pb_t^q)\bigr].
\end{equation}
This is a matrix polynomial identity, not a positivity approximation.


\paragraph{A version with one numerical matrix envelope.}
The same dictionary also applies on the more elementary sufficient domain
\[
 n(E_0)\le\tfrac1{12},\qquad n(C_{pq})\le128\quad((p,q)\in I).
\]
Indeed, put $T_\delta=2\delta+2\delta^2+\delta^3+\delta^4+\delta^5$.
At $\delta=2^{-12}$, rational arithmetic gives
$128T_\delta<1/12$. Replace the source matrices above by
\[
 H_0=C_{00}-128T_\delta I,\qquad
 H_{pq}^{\pm}=\tfrac12\delta^{p+q}(128I\pm C_{pq}).
\]
They are positive semidefinite, with $H_0\succeq I/6$, and the
same matrix identity holds. Here even the matrix envelope is a fixed
numerical constant. Section~\ref{sec:sourcebox} establishes these inequalities throughout
the printed source box. They are numerical source estimates, not
a consequence of pointwise boundedness alone.

\paragraph{Only two gate radicals.}
Write $a=A_t^4$, $b_t=B_t^4$, and introduce
\[
 \xi^2=1-A_t^4,\qquad \eta^2=1-B_t^4.
\]
For every nonnegative integer pair $(p,q)$,
\begin{align*}
 1+A_t^{4p}B_t^{4q}&=1+(A_t^{2p}B_t^{2q})^2,\\
 1-A_t^{4p}B_t^{4q}
  &=\sum_{i=0}^{p-1}(\xi A_t^{2i})^2
    +\sum_{j=0}^{q-1}(\eta A_t^{2p}B_t^{2j})^2.
\end{align*}
Empty sums mean zero. The second identity is a finite geometric sum.
For the seven pairs in $I$, the resulting common multiplier list is
\[
\begin{split}
\mathcal M=\{&1,B_t^2,B_t^4,A_t^2,A_t^4,A_t^4B_t^2,A_t^8,A_t^8B_t^2,\\
             &\eta,\eta B_t^2,\eta A_t^4,\eta A_t^8,
               \xi,\xi A_t^2,\xi A_t^4,\xi A_t^6\}.
\end{split}
\]
In particular there is no new radical for each monomial, and no
$\sqrt{M\pm y(t)}$ depending on an adaptively chosen $M$.

\paragraph{A compact algebraic cover of the testing chart.}
Set
\[
 u=\delta^2A_t^8,\quad w=\delta B_t^4,\quad W=\omega^4,\quad Z=\zeta^4,
 \qquad \omega^8+\zeta^8+X^2+Y^2=1.
\]
Together with the two gate equations, these equations define a compact
real algebraic testing set. They cover exactly the normalized real-state
chart $W,Z\ge0$, $0\le\sqrt u,w\le\delta$. All signs of
$A_t,B_t,\omega,\zeta,\xi,\eta$ are allowed. The common state sign takes
the same-sign state chart $WZ\ge0$ to this chart. In particular
\[
 \sqrt{uwWZ}=\delta^{3/2}A_t^4B_t^2\omega^2\zeta^2
\]
is the nonnegative root for every real point of the cover. Since
$\delta=2^{-12}$, every power $\delta^{j/4}$ appearing in the
substitution is rational.

A primitive rescaling of the twelve testing functions is
\[
\begin{split}
 \chi=(&X,Y,\zeta^4,\omega^4,A_t^8\zeta^4,B_t^4\zeta^4,
 A_t^8\omega^4,A_t^8B_t^4\omega^4,A_t^{16}\omega^4,\\
 &A_t^{12}\omega^4,A_t^4\omega^4,A_t^4B_t^2\omega^2\zeta^2).
\end{split}
\]
The circuit coordinates from the preceding section satisfy
$u^{1/4}=\delta^{1/2}A_t^2$ and $w^{1/4}=\delta^{1/4}|B_t|$.
Their squared identities remain valid when $|B_t|$ is replaced by
$B_t$: every monomial in each root has the same parity in $B_t$.
The same parity argument applies to $\omega$ and $\zeta$.

The original twelve-root vector is $g=R(c)\chi$ for the printed
source matrix, after multiplying its columns by
\[
 (1,1,1,1,\delta^2,\delta,\delta^2,\delta^3,
                 \delta^4,\delta^3,\delta,\delta^{3/2}).
\]
On the negative-state chart use $Z=-\zeta^4$, delete the final
coordinate of $\chi$, and take the first eleven columns of the
negative root matrix $A_-$ from \texttt{C1\_ROOT\_MATRIX.json}.
The exact source-to-feature map is
\[
 R_-=A_-[:,1{:}11]\,
 \diag(1,1,-1,1,-\delta^2,-\delta,\delta^2,\delta^3,
                                  \delta^4,\delta^3,\delta).
\]
Here $[:,1{:}11]$ means the first eleven columns in one-based
notation, equivalently \texttt{A\_minus[:, :11]} in Python.
These three negative column signs are necessary because the
feature vector uses $\zeta^4$, whereas the source root formulas
use the signed state coordinate $Z$.

Combining repeated multipliers in \eqref{eq:fixed-matrix-split} yields
\[
 g^{\mathsf T}Qg
   =\sum_{m\in\mathcal M}(m\chi)^{\mathsf T}G_m(c)(m\chi),
 \qquad G_m(c)\succeq0.
\]
Each $G_m$ is a sum of congruences $R(c)^{\mathsf T}H R(c)$ of the
explicit source matrices $H_0,H_{pq}^{\pm}$. Source-field square
completion therefore gives a finite separated SOS in this one fixed
testing family. Both state normals are preserved because the whole
matrix $Q$ is used, rather than a determinant or a selected restriction.
For the normalized physical slack, multiply every testing function by
$\kappa_t$, with
\[
 \kappa_t^2(1+\delta^4A_t^{16})(1+\delta^2B_t^8)=1.
\]
The denominator is strictly positive on the entire testing set.
This is a pure testing restoration, and the sign of $\kappa_t$ has
no effect on the Gram identity.

\paragraph{Evaluated rational matrices.}
The source files print sixteen feature blocks of length twelve, hence
192 feature slots. Exactly 152 polynomial functions are distinct.
Using only the gate relations, the block products span 280 independent
monomials. No state-sphere reduction is needed: every product is
homogeneous of state degree two, so a relation on the normalized
state chart would extend by radial scaling to an open cone and hence
would already be a polynomial relation. There are 1248 symmetric Gram entries. The fully evaluated
coefficient map has 1872 nonzero integer entries and exact row rank
280. Rational elimination produces an affine pencil of sixteen
$12\times12$ blocks in 280 coefficient variables and 968 auxiliary
variables. It has 2854 nonzero affine entries, with denominators only
one or two. Direct multiplication verifies all 280 coefficient equations.

For precision, this evaluated pencil represents the principal Gram
cone exactly: a coefficient vector is in that cone if and only if an
auxiliary vector makes all sixteen blocks positive semidefinite.
It is not a representation of the full quantum set. The retained
circuit, event, and angular roots belong to additional fixed blocks,
and a global extraction also requires source and testing coverage.

\paragraph{The retained circuit scalar uses the same gates.}
The scalar multiplying the nonnegative circuit remainder is
\[
 q_C=q_{C,0}-r_Cu,\qquad
 q_{C,0}=\tfrac12-\frac{c_Ad_G}{S}\left(\frac2\lambda+k\right),
 \quad r_C=\frac{c_Ad_G}{S}\ge0.
\]
The circuit diagonal of the principal Gram and
\eqref{eq:finite-admission} give $q_{C,0}-r_C\delta^2\ge1/3$.
Exactly,
\[
 q_C=(q_{C,0}-r_C\delta^2)
             +r_C\delta^2\bigl[\xi^2+(\xi A_t^2)^2\bigr].
\]
Thus multiplying the three elementary circuit roots requires only the
fixed multipliers $1,\xi,\xi A_t^2$; no adaptive radical is needed here
either.


\section{The evaluated full-parent local source pencils}
The preceding fixed multiplier construction was applied to every
principal term, and all retained roots from the complete C1 identity
were then restored as their own source-positive Gram blocks.
Each elementary or angular retained root has source/testing rank at
most three. The circuit roots are multiplied by the three fixed
functions $1,\xi,\xi A_t^2$ just derived. No guessed monomial extension
is used in this compilation.

For $Z\ge0$, the resulting fixed cone has 52 blocks, total matrix
order 268, maximum block order 12, and 168 distinct testing functions.
Its coefficient map has 281 independent rows and 1368 Gram entries.
Substituting the full parent, linear in the seven canonical source
coefficients $(d,m,C,A,D,\alpha,\beta)$, and eliminating the
coefficient equations gives an evaluated rational pencil with 1087
auxiliary variables. For $Z\le0$, the common circuit-centre term is
absent from the principal root matrix, so its testing span has eleven
functions. The corresponding fixed cone has 49 blocks, total order
240, maximum block order 11, 143 distinct functions, 221 independent
coefficient equations, and 934 auxiliaries.

Block-diagonal concatenation gives one evaluated rational pencil of
101 blocks and total order 508, in the same seven source coefficients
and 2021 auxiliaries. Every one of its 5326 nonzero upper-triangle
affine entries is printed. Direct rational multiplication verifies all
502 full-parent coefficient equations after elimination.

The equivalence proved by these matrices is exact and limited in
scope: a canonical source belongs to both displayed fixed Gram cones
if and only if this pencil admits a positive-semidefinite completion.
Feasible Grams imply nonnegativity of its full slack on both testing
state charts. Conversely, the uniform source estimates together with
the explicit source Gram assignments prove membership for every
admitted native source. The exact gauges
\[
 (u,Y,Z)\mapsto(-u,-Y,-Z),\qquad
 (w,X,Z)\mapsto(-w,-X,-Z)
\]
preserve the full parent and extend the conclusion to
$|u|\le2^{-24}$, $|w|\le2^{-12}$ and every real state. Applying the
real quadratic identity to real and imaginary state parts covers
complex states. The final polynomial identities specialize directly
to testing axes and zero state coordinates. No state coordinate,
$D$, or $\beta$ is inverted in the final family.

These two pencils are full-state local certificates. They do not
represent the full quantum set, and they do not establish that an
arbitrary globally valid C1 or C2 source belongs to the same fixed
cone. The uncovered source and testing regions remain a separate
constructive obligation.

